% Chapter 6 worked example. Included by ch06_resources.tex.
\section{Worked Example: Clinical Documentation Platform: Lifecycle Cost and Pricing}
\label{sec:comp-ch06-worked-example}

\subsection{Purpose and Status}
This worked example demonstrates how the methods in Chapter~6 can be
translated into a reviewable institutional decision. All numerical inputs are
synthetic unless a source is identified explicitly. They are intended to show the
logic of the analysis, not to report empirical findings or establish a universal
threshold. Local use requires replacement of the assumptions with current,
versioned evidence and approval through the relevant clinical, managerial,
economic, legal, technical, and governance processes.

A hospital group is comparing three pricing proposals for a documentation platform: a fixed enterprise licence, a per-clinician subscription, and a hybrid base fee plus per-encounter charge. Each proposal excludes some implementation, monitoring, and exit activities. 

\subsection{Decision Problem and Comparators}
\textbf{Decision problem.} Which pricing and contracting structure produces the most predictable and defensible lifecycle cost, and under what conditions could projected benefits be realized?

The comparator set is deliberately broader than the existing pathway. This prevents
a new technology from receiving a lower evidentiary threshold than staffing,
workflow, shared-service, conventional digital, or no-change alternatives.

\begin{longtable}{@{}p{0.08\textwidth}p{0.84\textwidth}@{}}
\caption{Comparators for the Chapter 6 worked example}
\label{tab:comp-ch06-comparators}\\
\toprule
\textbf{Option} & \textbf{Comparator or strategy} \\
\midrule
\endfirsthead
\toprule
\textbf{Option} & \textbf{Comparator or strategy} \\
\midrule
\endhead
\bottomrule
\endlastfoot
1 & Fixed enterprise licence \\
2 & Per-clinician subscription \\
3 & Hybrid base plus encounter-based fee \\
4 & Do not procure; invest in workflow and template redesign \\
\end{longtable}

\subsection{Model and Decision Structure}
The analytical structure should follow the decision rather than the available
software. The team should map the causal pathway from intervention input to
information, human decision, action, outcome, resource consequence, and review.
The structure should also identify capacity constraints, uncertainty, subgroup
variation, implementation requirements, and events that would change the
recommendation.

A compact quantitative relationship used in this example is:
\begin{equation}
\mathrm{ROI}=\frac{\mathrm{PV}(B)-\mathrm{PV}(C)}{\mathrm{PV}(C)}\times100\%
\label{eq:comp-ch06-key-relationship}
\end{equation}
The equation is an organizing aid rather than a substitute for disaggregated evidence,
qualitative findings, or mandatory safeguards. Its variables and interpretation must
be defined in the local protocol.

\subsection{Synthetic Parameters and Evidence Sources}
Table~\ref{tab:comp-ch06-parameters} provides the base-case inputs. A
production analysis should add parameter distributions, correlations, structural
alternatives, data dates, system versions, and evidence-confidence ratings.

\begin{longtable}{@{}p{0.32\textwidth}p{0.22\textwidth}p{0.38\textwidth}@{}}
\caption{Synthetic parameters for the Chapter 6 worked example}
\label{tab:comp-ch06-parameters}\\
\toprule
\textbf{Parameter} & \textbf{Base value} & \textbf{Source or interpretation} \\
\midrule
\endfirsthead
\toprule
\textbf{Parameter} & \textbf{Base value} & \textbf{Source or interpretation} \\
\midrule
\endhead
\bottomrule
\endlastfoot
Clinicians & 420 & Workforce plan \\
Annual eligible encounters & 520,000 & Activity forecast \\
Enterprise base fee & EUR 1.15 million/year & Vendor A \\
Per-clinician fee & EUR 390/month & Vendor B \\
Hybrid base fee & EUR 540,000/year & Vendor C \\
Hybrid variable fee & EUR 1.15/encounter & Vendor C \\
Initial integration and governance & EUR 760,000 & Local estimate \\
Annual monitoring and support & EUR 310,000 & Local estimate \\
Expected useful decision horizon & 5 years & Business case \\
\end{longtable}

\subsection{Step-by-Step Analysis}
The sequence below is intended to make the analysis reproducible and to ensure
that evidence, economics, governance, and implementation develop together.

\begin{longtable}{@{}p{0.07\textwidth}p{0.55\textwidth}p{0.30\textwidth}@{}}
\caption{Analytical sequence}
\label{tab:comp-ch06-analysis-steps}\\
\toprule
\textbf{Step} & \textbf{Required work} & \textbf{Decision record} \\
\midrule
\endfirsthead
\toprule
\textbf{Step} & \textbf{Required work} & \textbf{Decision record} \\
\midrule
\endhead
\bottomrule
\endlastfoot
1 & Define perspective, price year, time horizon, comparator, and full cost boundary. & Evidence and responsible owner to be recorded \\
2 & Classify costs as fixed, step-fixed, variable, contingent, shared, or exit related. & Evidence and responsible owner to be recorded \\
3 & Model low, expected, and high activity for each pricing structure. & Evidence and responsible owner to be recorded \\
4 & Discount costs and credible monetised benefits; preserve non-monetised value separately. & Evidence and responsible owner to be recorded \\
5 & Create a benefit-realization register that distinguishes cash savings, avoided costs, usable capacity, quality, workforce, and public value. & Evidence and responsible owner to be recorded \\
6 & Test price, uptake, review burden, implementation delay, support, and exit thresholds. & Evidence and responsible owner to be recorded \\
\end{longtable}

\subsection{Illustrative Outputs}
The outputs in Table~\ref{tab:comp-ch06-outputs} show how technical,
clinical, operational, economic, equity, and governance findings should be kept
visible rather than compressed into a single favourable or unfavourable score.

\begin{longtable}{@{}p{0.30\textwidth}p{0.24\textwidth}p{0.38\textwidth}@{}}
\caption{Illustrative model and decision outputs}
\label{tab:comp-ch06-outputs}\\
\toprule
\textbf{Output} & \textbf{Result} & \textbf{Interpretation} \\
\midrule
\endfirsthead
\toprule
\textbf{Output} & \textbf{Result} & \textbf{Interpretation} \\
\midrule
\endhead
\bottomrule
\endlastfoot
Most predictable annual spend & Enterprise licence & May overpay at low uptake \\
Lowest expected five-year PV cost & Hybrid proposal & Dependent on activity cap and included support \\
Dominant uncertainty & Net usable capacity per encounter & Not vendor price alone \\
Cash-releasing saving & Not established & Released time is not yet linked to budget reduction \\
Preferred decision & Negotiate hybrid with cap and conditional renewal & Preserves predictability and proportionality \\
\end{longtable}

\subsection{Sensitivity, Uncertainty, and Subgroups}
At minimum, the completed analysis should vary the parameters most likely to
change the decision, test at least one credible alternative model structure, and
report subgroup results separately where performance, access, response, burden,
or opportunity cost may differ. Deterministic analysis should identify thresholds;
probabilistic analysis should be used where credible parameter distributions and
correlations are available. Scenario analysis is preferable when structural or
future uncertainty cannot be represented defensibly by one probability model.

The record should distinguish uncertainty that can be reduced through evidence,
uncertainty that can be managed through safeguards, and uncertainty that remains
too consequential to accept. Missing evidence should not be represented as an
average or neutral value without justification.

\subsection{Interpretation and Recommendation}
Negotiate a capped hybrid structure with explicit integration, monitoring, updates, data access, and exit support. Treat productivity benefit as conditional until local evidence shows how net time is converted into usable capacity or another measurable outcome.

The recommendation is conditional rather than permanent. It applies only to the
specified population, setting, intervention configuration, evidence cut-off, and
version. The following conditions should be incorporated into the approval,
contract, implementation, monitoring, or renewal record:

\begin{itemize}
 \item all pricing units are clinically and financially understandable.
 \item base and variable components do not double charge essential services.
 \item low, expected, and high activity scenarios remain affordable.
 \item benefit owners and verification sources are named.
 \item renewal depends on observed value, current evidence, and acceptable exit options.
\end{itemize}

\subsection{Decision Statement}
\begin{quote}
For the specified decision problem and comparator set, the institution should
\emph{[proceed / proceed with conditions / redesign / pilot / restrict / defer /
replace / retire]}. The decision applies to \emph{[population, setting, version,
and evidence period]}, is owned by \emph{[accountable authority]}, and will be
reviewed on \emph{[date]} or earlier if \emph{[trigger events]} occur. The
evidence, assumptions, unresolved uncertainty, mandatory safeguards, monitoring
thresholds, and exit arrangements are recorded in the accompanying dossier.
\end{quote}
